\section{Related Work}
\label{app:related}

\paragraph{Prototype and memory methods.} Prototype and memory methods represent data by stored exemplars or centers. DGM uses the same computational intuition but adds a refinement rule: a new prototype is created when local repair cannot absorb evidence without conflating consequences, and distinction edges record the separating evidence.

\paragraph{Decision trees and splitting criteria.} Classical decision trees use information gain and variance reduction as splitting criteria. In our view, a decision tree is a hard-routing special case of a distinction graph. It stores distinctions in a hierarchy rather than as local edges between concepts.

\paragraph{Sparse experts and retrieval layers.} Mixture-of-experts and retrieval-augmented models scale neural computation by selecting a small subset of parameters or memories for each input. A DGM layer has a similar sparse read path, but the selected objects are runtime concepts with local memories, and the write path can refine or repair the state after observing consequences.

\paragraph{Information theory and statistical decision theory.} The log-loss identity is the standard relation between optimal code length and conditional entropy, while the squared-loss identity is the orthogonal decomposition of conditional expectations in $L^2$. Our contribution is to place these identities inside a runtime-learning story: the value of learning a distinction is its reduction in Bayes risk.

\paragraph{Backpropagation and gradient learning.} Backpropagation supplies differentiable credit assignment for fixed computation graphs \citep{rumelhart1986learning}. Distinction Learning addresses a different failure mode: the current information structure may not contain the distinctions needed for any measurable predictor to succeed.

\paragraph{Projection learning and online repair.} Bregman projection, mirror descent, and proximal methods provide minimal-change state updates \citep{bregman1967relaxation,zinkevich2003online,shalev2012online,hazan2016introduction}. Conservative Projection Memory uses this idea for runtime associative repair. The present paper demotes projection repair to the continuous-state layer that operates after the relevant information structure has been fixed.

\paragraph{Test-time adaptation.} Test-time training adapts parameters using auxiliary losses at inference time \citep{sun2020test}. TTT layers treat hidden states as trainable models updated on test sequences \citep{sun2025learning}. Distinction Learning suggests a complementary diagnostic: when runtime evidence conflicts inside an indistinguishable atom, the issue is not only how to update fast state, but what distinction the state failed to represent.
